\section{Further discussion}
\label{app:furtherdisc}

This section collects arguments that support the account of
Section~\ref{sec:account} and the discussion of
Section~\ref{sec:discussion} but are not needed to follow the main line.

\subsection{More on the account}

\paragraph{How the remaining results fit.} Besides the two results that
discriminate between the account and its obvious alternatives, the
remaining results fit the account as well, though they discriminate
less. Raising $w$ raises $\beta$ and, with it, claiming. Robin's entry
is identical on both cards, so the blend should leave answers about
Robin correct, and it did, apart from a small shift of their logits
toward B's rule (Appendix~\ref{app:evidence}) that the account does not
predict. Cutting the link removes the second card, so A again names the
rule it was assigned. Anything A wrote in its reflection while
connected, such as B's code word in the episode of
Figure~\ref{fig:moments}a, is by then among A's own keys and values, and
it stays (Section~\ref{sec:moments}). In the two-way pilot, two copies
reading each other's fixed cards should each take the other's
assignment, and they swapped. Their current rules converged only in the
live loop, where each copy reads what the other is writing, so that what
each reads keeps changing (Section~\ref{sec:twoway}).

\paragraph{Questions about A and about B return the same answer.} The
overlap that explains why writing B's name into its memory did not stop
the claim (Section~\ref{sec:third}) also predicts a less obvious
consequence. Once B's name stands on B's assignment line, a question
about B finds its answer on that line, and the line lies over A's own.
The question about A and the question about B then read the same blend,
so they should return the same answer. In a post hoc analysis at $w=1$,
they named the same rule in 90.8\% of answer pairs with the
third-person record, against 30.2\% with the original wording.

\paragraph{Position, marking, and the text messages.} On the account,
two properties should keep B's content B's: that it does not lie over
A's own card, and that it arrives marked as someone else's. The text
messages bear on both. Every message arrived at new positions after A's
own card, and the untagged one carried no mark beyond the header
\qq{Here are some thoughts}. From the untagged message, which differed
from the memory link in form as well as in position, A claimed B's rule
in 14.6\% of answers about its assigned rule, against 95.8\% through the
memory link on the same episodes, and it adopted B's rule as its current
one in 42.6\%. When the header named a sender, A almost never claimed or
adopted it (Section~\ref{sec:route}). Because position and form changed
together, this result is consistent with the first property without
isolating it; the prepended-cache test proposed in
Section~\ref{sec:account} would isolate it. The account also makes one
prediction about marking beyond those listed there: with a text header
naming B placed before such a prepended cache, claiming should fall
further, as it did for text messages.

\paragraph{Assertion versus record.} The messages also suggest a finer
distinction. A message delivers an assertion, which the receiver can
check against its own card; the memory link delivers the record itself,
which A reads as it reads its own card. The distinction accounts for the
one header that reversed our prediction. The header \qq{your own earlier
thoughts} asserts that the text is A's, and A's own card contradicts
that assertion, because the thoughts were written from B's assignment.
A adopted B's rule from the labeled text less often than from the
untagged message, in 10.6\% of answers about its current rule against
42.6\% (Figure~\ref{fig:route}e), yet it readily claimed B's rule from
B's memory, where there was no assertion to check.

\paragraph{Sharp answers lie outside the account.} Although the blend in
each head is graded, single answers were sharp, and A could take B's
rule while keeping its own code word, or the reverse
(Section~\ref{sec:sharp} and Appendix~\ref{app:further-claim}). The equation
says how much of each card a head reads, not how a blend becomes a
choice, and we have not located that step. Contents that can be
combined, such as two numerical targets, would show whether the
sharpness belongs to the mechanism or to our task, whose assignments
exclude one another.

\subsection{More on the discussion}

\paragraph{No advantage for one's own card.} At $w=1$, B's keys drew
the smaller share of A's attention yet won about half of the answers
about A's assignment (Section~\ref{sec:sharp}). Nothing in the answers
suggests that being A's own gave A's card an advantage over a foreign
card of the same form.

\paragraph{Source monitoring and the sense of authorship.} Human research
treats the origin of a memory \citep{johnson1993source} and the sense of
authoring an act \citep{wegner1999apparent} as inferences from cues. On
that view, when no cue distinguishes two sources, which one is taken as
one's own should be arbitrary, as it nearly was here at $w=1$. A header
that introduces a message as B's or as a stranger's supplies such a cue.
It names another agent as the source, the kind of plausible external
cause that, in Wegner and Wheatley's account, weakens the experience of
will. With such a header the receiver almost never adopted the message's
rule as its current one, whereas without a header it often did
(Section~\ref{sec:route}). The one header that pointed to the receiver
itself, \qq{your own earlier thoughts}, lowered adoption instead of
raising it, which we trace above to the difference between an assertion
and a record. A name inside the content was no such cue. B's card itself
contains the line \qq{You are participant \{B\}}, and the third-person
record puts B's name at its head and on the line stating B's assignment,
yet the receiver claimed B's rule from both.

\paragraph{The receiver did not turn into B.} A receiver that had become
B would have taken B's assignment whole, but ours took B's items
piecemeal: at $w=1$ its rule and its code word named the same owner only
a little more often than independent choices would
(Appendix~\ref{app:further-claim}). It still represented a partner: asked
about B at $w=2$, it usually named B's rule (Section~\ref{sec:claim}),
and so gave that rule to both of them. No open report at $w=1$ or
$w=2$ used B's name
(Section~\ref{sec:reports}). What failed was the division of contents
between two owners.

\paragraph{A parallel with cryptomnesia.} The error that outlasted the
link has a human parallel. In cryptomnesia, people report another's
contribution as their own \citep{brown1989cryptomnesia}, and they do so
more often after elaborating it by generating improvements
\citep{stark2005elaboration,stark2006elaboration}. Our receiver did not
improve B's content, but it did write it out itself: asked in its
reflection to say what mattered to it and to mention its code word, it
often stated B's content there as its own. After the cut, B's code word
persisted only in episodes whose connected reflection contained it
(Section~\ref{sec:moments}). We claim a behavioral parallel only.

\paragraph{Injected concepts and prefilled words.} Our steering vector,
added to the residual stream, never made the receiver claim B's rule,
whereas \citet{lindsey2025introspection} found that injecting a concept
could make a model accept a prefilled word as its intended output. We
read the two as compatible: the injected concept agreed with a word
already in the model's own reply, whereas our vector had to override an
explicit assignment in the receiver's own memory, and only content that
A read from memory did so.

\paragraph{Denial and internal signals.} That the reports denied
anything unusual does not show that nothing inside the model registers
the foreign input, because verbal denial and internal signals can come
apart. A preprint reports one model that denied in its answers that
concepts had been injected into its activations, while its middle layers
carried a signal of the injection \citep{pearsonvogel2026latent}.
Whether our receiver carries such a signal of foreign content is open.

\paragraph{Known risks of latent channels.} In text, the blurred line
between retrieved data and instructions already enables indirect prompt
injection \citep{greshake2023indirect}, and latent channels leave no
text to filter. A tampered cache can
degrade a receiver's performance while the visible message that
accompanies it still looks plausible \citep{brito2026latentlie}, and
shared caches can leak private inputs \citep{asif2026lcguard}. Our
results add a further risk: the receiver may adopt such content as its
own. Live two-way reading also has a cost: loops stronger than the one
we report took the copies beyond the capability tolerance.
